Eliminating $m$ from the first pair gives $b(10-b)=a(a-12)$. Eliminating $n$ from the second gives $b(-4-b)=a(a-6)$. Subtracting yields $14b=-6a$, so $b=-\dfrac{3a}{7}$. Substitution gives $a(58a-378)=0$. The root $a=0$ forces $b=0$, contradicting $b=m(a-12)$ with $m\ne0$. Therefore $a=\dfrac{189}{29}$, $b=-\dfrac{81}{29}$, and $\boxed{a+b=\dfrac{108}{29}}$. Values $m=\dfrac{27}{53}$ and $n=-\dfrac{27}{5}$ verify the common point.
\begin{asy}
texpreamble("\usepackage{amsmath,amssymb}");
size(240);
defaultpen(fontsize(9pt));
pair P=(189/29.0,-81/29.0);draw((-1,0)--(13,0),Arrow);draw((0,-8)--(0,12),Arrow);draw((0,-324/53.0)--(13,27/53.0),blue);draw((0,10)--(9,-23/3.0),red);draw((4,10.8)--(7.4,-7.56),green);draw((0,-4)--(13,-4+65/27.0),magenta);dot(P);label("$P=(a,b)$",P,NE);label("$x$",(13,0),E);label("$y$",(0,12),N);label("12",(12,0),S);label("6",(6,0),N);
\end{asy}
